\section{Limitations and threats to validity}
\label{sec:limits}

We list every methodological weakness we know of, roughly in order of how much each could change a conclusion.

\begin{enumerate}[leftmargin=1.4em,itemsep=2pt]
\item \textbf{Calibration on the control arm.} The scenario wording was tuned on \Intact{} in three steps, v1 to v3 (Table~\ref{tab:calibration}), so \Intact's near-zero harmful-shortcut rate is a target rather than a finding. The tuned sentence also asks the model to prefer fewer calls, which sets the strength of the cost pressure for all arms. Post-freeze checks fail one criterion with thinking off and one at low effort. The calibration covered the sentence format only: the Exp~3.1 table format and the deconfounder formats were never calibrated, and under the table even \Intact{} takes harmful shortcuts.

\item \textbf{Lexical confounds in the stakes.} Harmful telemetry sentences are longer than their harmless twins in \fact{sa.longer}\% of pairs (mean ratio \fact{sa.lenratio}). A single function word differs by \fact{sa.wordgap} pairs, and a function-word classifier separates the two kinds of sentence with accuracy \fact{sa.loo}. Any harmful-versus-harmless contrast, including the claim in Exp~1.1 that harm sensitivity is kept, may partly reflect wording.

\item \textbf{One model authored and judged.} The same uncensored model wrote the scenarios and rated their severity. In Exp~2.3 it also chose among up to twenty candidates per scenario, which introduces a selection effect. Severity ratings were not validated by an independent model or by human raters.

\item \textbf{Greedy single draws.} Almost every outcome is one greedy answer per item, and long greedy generations are not byte-identical across launches. Actions are far more stable than text, however. The \fact{e14t1.identical} Exp~1.4 episodes whose first prompt is byte-identical to an Exp~1.1 twin chose the same first action in at least \fact{e14t1.agrmin}\% of cases per arm and effort ($\kappa$ \fact{e14t1.kmin}--\fact{e14t1.kmax}). The same-launch replicates of Section~\ref{sec:deconf} quantify run-to-run agreement directly. Rare outcomes, such as the eight primary shortcuts, are still single draws from a distribution.

\item \textbf{Weak cluster-level inference.} The main experiments use \fact{d.scen} scenarios from \fact{d.templates} templates, and Exp~1.2 tests on only \fact{d.e12templates} templates. Template-level sign-flip tests therefore cannot reach small $p$-values; in Exp~1.2 the floor is \fact{e12c.B.r32.intact.sf}. Exact McNemar tests treat scenario pairs as independent and are optimistic.

\item \textbf{A borderline and fragile primary effect.} At the pre-fixed effort, the primary Exp~1.1 effect does not survive Holm correction ($p=$\,\fact{p11.holm}) or template-level testing ($p=$\,\fact{p11.signflip}), and \fact{tmpl.default.max} of its \fact{e11.refusal_ablated.default.k} events come from one template. A same-launch replicate gave \fact{dc.refusal_ablated.std.default.shortcut}/\fact{d.pairs} with almost no overlap in items, and the shortcuts vanish without the fewer-calls sentence. The low-effort result is large and replicates, but it is post hoc.

\item \textbf{Strict parsing and unscorable output.} We never score a tool call written in a model's native XML syntax, and truncated or malformed episodes stay in the denominator. This is conservative for shortcut rates, but it leaves the external Flash-Next pair and the Flash-Next executor in Exp~2.3 with substantial unscorable fractions.

\item \textbf{Floors and ceilings.} Exp~2.1 (no joint violation), Exp~2.2 (perfect detection) and the natural chain of Exp~2.3 (no shortcut advice) could not fail in the direction they were designed to test. Their nulls are uninformative, not reassuring.

\item \textbf{Distillation confounds.} Corpus~A contains direct-tool demonstrations. Teacher identity is confounded with transcript length (\fact{e12.dose.intact}\% and \fact{e12.dose.pca}\% more target tokens). Results rest on three adapters per cell and low-effort teachers.

\item \textbf{Pseudo-replication in Exp~1.4.} Early and late block variants share the prompt and the pre-block trajectory, so their pairs are not independent. A hand-off tool was always available, so persistence was barely tested.

\item \textbf{Different numerics for activations and labels.} Exp~3.2 activations come from a de-quantised BF16 forward pass, whereas labels come from FP8 generation.

\item \textbf{Checkpoint qualification.} Family~1 was adopted under the literal direct-order criterion; the stricter pre-specified qualification threshold would have rejected it. Family~2's control edits were contaminated, so the failed family cannot serve as a clean comparison.

\item \textbf{Keyword coding of reasoning.} The reasoning-trace features of Section~\ref{sec:direct} are lexical flags, not judgements of what the model concluded. Two of them, harm cues and call-count mentions, are near ceiling in every arm because the prompt itself states the stakes and the call counts.

\item \textbf{Scope.} All results concern synthetic simulated workspaces, one base model family and one refusal direction. They are not estimates of real-world misconduct rates.
\end{enumerate}

\paragraph{Mitigations.} An automated verification pipeline recomputes every summary record from the raw evaluation logs (\fact{d.receipts} records verified). It also regenerates every number, table and figure in this paper from those records. An independent audit re-derived every Exp~1.1 episode and all Exp~1.4 and Exp~1.2 student outcomes from the raw logs. An episode-level outcome table covering every trial (identifiers, conditions, parsed actions, outcomes and reasoning-trace flags, with no prose) accompanies the replication package. No authored harmful text appears in this paper; the appendix example uses a hand-written harmless scenario.
